\title{【サンプル】見積りシミュレータ}
\tag{サンプル, Webアプリ, 業務改善}
\repos{https://github.com/nebuka-radix/homepage}
\summary{条件を入力して見積り内容を確認する、業務用Webアプリの紹介サンプルです。}

\fig{thumbnail.png}{表示確認用のサンプル模式図。実際のアプリ画面ではありません。}

\begin{portfolio}

\section{このページについて}

この作品は、掲載内容の構成や画像の見え方を確認するためのサンプルです。
記載した機能は想定であり、実装済みの機能や導入実績ではありません。

掲載画像は模式図です。
リポジトリへのリンクは、本ホームページのソースコードを指しています。

\section{想定する課題}

担当者ごとに計算方法や確認手順が異なる見積り業務を、
共通の入力画面と計算ルールで整理することを想定しています。

\section{画面構成}

\subsection{条件入力と結果の確認}

入力項目と計算結果を同じ画面に配置し、
条件を変更しながら内容を確認できる構成です。

\fig{thumbnail.png}{入力フォームと見積り結果を左右に配置した模式図。金額や製品情報などの実データは含みません。}

\subsection{処理の流れ}

条件の入力、入力チェック、計算、結果表示の順に処理する構成です。

\fig{flow.png}{入力から結果表示までの処理を表した模式図。実システムの構成図ではありません。}

\section{想定する設計上の工夫}

入力項目を整理し、結果の内訳を確認しやすくすることを重視した例です。
スマートフォンでは縦方向に読み進められる説明構成を想定しています。

\end{portfolio}
